\section{Реализация словаря}
\label{sec:implementation}
\begingroup
\raggedright

Для каждого ключа дерево хранит один узел со словом в нижнем регистре, значением, номером бита и двумя ссылками. При поиске бит ключа указывает, по какой из двух ссылок идти дальше. Как только номер бита перестаёт расти, спуск заканчивается. Затем программа сравнивает найденное слово с искомым целиком.

При вставке программа сравнивает новое слово с тем, к которому привёл поиск. Первый бит, в котором слова различаются, определяет место нового узла. При удалении узла с завершающей ссылкой на себя достаточно обойти его по другой ветви. Если завершающая ссылка пришла из другого узла, тот занимает место удаляемого.

Если длина всех ключей ограничена \(L\) байтами, то за спуск проверяется не более \(8L\) битов. Поиск, вставка и удаление требуют постоянного числа спусков и сравнений строк, поэтому занимают \(O(L)\) времени в худшем случае.

\clearpage
\subsection{Пример работы}

Словарь распознаёт \texttt{ALPHA} как уже добавленное слово и сохраняет значение 34. После удаления поиск возвращает \texttt{NoSuchWord}. Загрузка восстанавливает обе пары из файла. Программа собирает отдельное дерево и только после успешной загрузки заменяет им рабочее.

\begin{minipage}[t]{0.53\linewidth}
\textbf{Вход}
\begin{lstlisting}[language={},numbers=none]
+ Alpha 34
+ ALPHA 99
alpha
+ Alphabet 0
! Save dictionary.bin
- alpha
ALPHA
! Load dictionary.bin
alpha
ALPHABET
\end{lstlisting}
\end{minipage}\hfill
\begin{minipage}[t]{0.40\linewidth}
\textbf{Вывод}
\begin{lstlisting}[language={},numbers=none]
OK
Exist
OK: 34
OK
OK
OK
NoSuchWord
OK
OK: 34
OK: 0
\end{lstlisting}
\end{minipage}
\par
\endgroup
